\chapter{Indices de crédit et tranches}\label{ch:m2:credit-indices-and-tranches}

Au premier trimestre 2012, le Chief Investment Office de JPMorgan Chase a triplé
le notionnel net d'un portefeuille d'\omterm{def:m2:credit-indices-and-tranches:index}{indices de crédit} et de \omterm{def:m2:credit-indices-and-tranches:tranche}{tranches}, créé en 2007 comme
couverture contre les pertes de crédit, pour atteindre environ 157~milliards de dollars. Une large part correspondait à une protection
vendue sur le contrat à dix ans d'une \omterm{def:m2:credit-indices-and-tranches:index}{série d'indice} nord-américain investment grade,
émise en 2007~; à un moment donné, la position nette sur cette série atteignait
82~milliards de dollars, soit 10 à 15 jours du volume négocié sur l'ensemble du marché.
Ses traders en vinrent à croire que le marché connaissait leurs positions et fixait ses prix
à leur détriment. La banque révéla une perte d'environ 2~milliards de dollars en mai~; à la fin de
l'année, les pertes atteignaient 6.2~milliards de dollars, et elle reconnut plus tard les faits dans un
accord amiable avec la SEC. Ce chapitre présente les instruments~: les \omterm{def:m2:credit-indices-and-tranches:index}{indices de crédit} et
leur roll, l'écart entre un indice et ses constituants, les \omterm{def:m2:credit-indices-and-tranches:tranche}{tranches} qui découpent les
pertes d'un indice, et les options sur indices.

\section{Les familles d'indices et le roll}

\begin{definition}[Indice de crédit, série d'indice]\label{def:m2:credit-indices-and-tranches:index}
Un \emph{indice de crédit}\index{indice de crédit} est un \omterm{def:m2:credit-default-swaps:cds}{credit default swap} (CDS) standardisé portant sur un
panier équipondéré d'\omterm{def:m2:credit-default-swaps:cds}{entités de référence}, avec un seul coupon et une seule maturité~; un
événement de crédit affectant un constituant est réglé sur la part de ce constituant dans le notionnel,
et l'indice se poursuit sur les autres. Une \emph{série d'indice}\index{série d'indice} est
une liste fixe de constituants, publiée à une date de roll~; une nouvelle série, dotée d'une
liste actualisée, est publiée tous les six mois, et la plus récente est la série on-the-run.
\end{definition}

Deux familles dominent~: CDX, indices nord-américains et de marchés émergents, et iTraxx,
européens et asiatiques, chacun en versions investment grade et high yield. L'indice
nord-américain investment grade regroupe 125 des noms les plus liquides de catégorie investment grade,
de poids égaux, et se renouvelle les 20 mars et 20 septembre~; un nom ayant
été déclassé, cessé d'être négocié ou fait l'objet d'une OPA est remplacé dans la nouvelle série,
tandis que l'ancienne série conserve sa liste et continue de se négocier, de moins en moins, en off-the-run.
L'indice verse un coupon fixe de 100 points de base par trimestre et il est coté en spread, avec
l'upfront réglé comme pour un nom individuel (chapitre~23). La série 2007 négociée par
le bureau de la banque en 2012 était la neuvième~; la série émise en septembre 2012 était la
dix-neuvième.

Un indice permet à un investisseur de négocier le crédit dans son ensemble~: une transaction sur un
instrument liquide achète ou vend de la protection sur un marché entier, bien moins cher que 125
noms individuels. Il peut servir à couvrir un book de noms individuels ou un portefeuille obligataire, ou à exprimer
une vue sur l'économie.

\begin{dated}{2026-09}{Caractéristiques des contrats d'indices}\label{dat:m2:credit-indices-and-tranches:indices}
CDX nord-américain investment grade~: 125 noms, poids égaux, renouvellement les 20 mars et
20 septembre, maturités de 1 à 10 ans, coupon de 100 points de base, compensé chez ICE Clear
Credit. Options sur indices~: de style européen, à règlement-livraison physique, sur le CDX à cinq ans
investment grade et high yield, iTraxx Europe et iTraxx Crossover, avec une échéance
pouvant aller jusqu'à neuf mois, compensées chez ICE.
\end{dated}

\section{Indice et constituants}

Un indice est un panier de CDS, et sa valeur se déduit donc des prix de ces CDS.
La cotation de l'indice elle-même, fixée par la négociation sur l'indice, n'a pas de raison de coïncider.

\begin{definition}[Spread intrinsèque, skew d'indice]\label{def:m2:credit-indices-and-tranches:intrinsic}
Le \emph{spread intrinsèque}\index{spread intrinsèque} d'un indice est le spread qui,
converti avec le coupon et les conventions de l'indice, donne la moyenne pondérée des
upfronts de ses constituants à ce coupon. Le \emph{skew d'indice}\index{skew d'indice} est le
spread coté de l'indice moins son spread intrinsèque.
\end{definition}

\begin{proposition}[Spread intrinsèque comme moyenne pondérée]\label{prop:m2:credit-indices-and-tranches:intrinsic}
Si le constituant $i$ a pour poids $w_i$, pour spread $s_i$ et pour annuité risquée $A_i$, l'indice
avec coupon $c$ a un upfront intrinsèque $\sum_i w_i (s_i - c) A_i$, et, au premier
ordre, le \omterm{def:m2:credit-indices-and-tranches:intrinsic}{spread intrinsèque}
\[
s^{\text{int}} \approx \frac{\sum_i w_i A_i s_i}{\sum_i w_i A_i},
\]
ce qui est inférieur à la moyenne simple des spreads, puisque les noms à spreads plus élevés ont
des annuités plus faibles.
\end{proposition}

\begin{proof}
L'upfront de chaque constituant vaut $(s_i - c)A_i$ selon le chapitre~23, et les jambes de
protection et de prime de l'indice sont les sommes pondérées de celles des constituants, donc son
upfront est leur somme pondérée. En l'écrivant $(s^{\text{int}} - c)A$ avec $A \approx
\sum_i w_i A_i$ donne la formule. Un spread plus élevé signifie une intensité de défaut plus élevée, et donc
une durée de vie espérée plus courte, d'où une annuité plus petite et un poids plus faible.
\end{proof}

\begin{example}[Un indice stylisé]\label{ex:m2:credit-indices-and-tranches:intrinsic}
Prenons 125 noms dont les spreads à cinq ans suivent une loi log-normale autour d'une médiane de 55
points de base, de 6.6 à 459, avec une moyenne simple de 75.3. Avec un taux plat de 4~\%
et un \omterm{def:m2:corporate-bonds:covenant}{taux de recouvrement} de 40~\%, leur upfront moyen au coupon de 100 points de base vaut
$-1.157\%$, ce qui donne un \omterm{def:m2:credit-indices-and-tranches:intrinsic}{spread intrinsèque} de 73.6 points de base~; la
moyenne des spreads pondérée par les annuités vaut aussi 73.6. Si l'indice se négocie à 65.6, son
skew vaut $-8$ points de base (voir la \cref{fig:m2:credit-indices-and-tranches:constituents}).
\end{example}

\begin{omfigure}
\begin{tikzpicture}
\begin{axis}[width=11cm,height=5cm,
  xlabel={constituant, trié par spread},ylabel={spread à cinq ans (pb)},
  tick label style={font=\small},label style={font=\small},
  xmin=0,xmax=126,ymin=0,ymax=300,grid=major,grid style={gray!25},
  legend columns=2,legend style={font=\footnotesize,at={(0.5,-0.3)},anchor=north,draw=none,fill=none,/tikz/every even column/.append style={column sep=8pt}},
  table/col sep=comma]
\addplot[omDef,only marks,mark=*,mark size=1pt] table[x=rank,y=spread]{figdata/markets-2/24-credit-indices-and-tranches/constituents.csv};
\addplot[gray,thick,dashed] coordinates {(0,75.3) (126,75.3)};
\addplot[omThm,very thick] coordinates {(0,73.6) (126,73.6)};
\addplot[omProp,very thick,dotted] coordinates {(0,65.6) (126,65.6)};
\legend{constituants,moyenne simple 75.3,intrinsèque 73.6,cotation de l'indice 65.6}
\end{axis}
\end{tikzpicture}

\omcaption{Les constituants d'un indice investment grade stylisé à 125 noms, triés par
spread à cinq ans (les deux plus larges, à 335 et 459 points de base, sont hors graphique). Le
\omterm{def:m2:credit-indices-and-tranches:intrinsic}{spread intrinsèque} est inférieur à la moyenne simple parce que les noms les plus larges portent
des annuités plus faibles~; un indice coté à 65.6 se négocie 8 points de base sous sa valeur intrinsèque.
Illustratif~; données~: le tutoriel du chapitre.}\label{fig:m2:credit-indices-and-tranches:constituents}
\end{omfigure}

Le skew est un arbitrage en principe~: acheter la protection sur l'indice et la vendre sur les
125 noms, et la position ne rapporte rien en cas de défaut, puisque chaque défaut est réglé
des deux côtés, et gagne le skew lorsqu'il se referme. En pratique, le skew persiste et
peut s'élargir. Négocier 125 noms individuels coûte bien plus qu'une seule transaction sur l'indice~; les noms individuels
sont beaucoup moins liquides que l'indice~; la transaction immobilise de la marge et du capital aussi longtemps
que le skew reste ouvert~; et un acheteur ou vendeur important et persistant sur l'indice peut le
faire s'écarter davantage avant qu'il ne revienne. Un skew est le point de rencontre du capital d'arbitrage et de la
demande d'un seul côté. Un très gros \omterm{def:m2:credit-default-swaps:cds}{vendeur de protection} sur indice fait tomber l'indice sous ses
constituants, et les investisseurs qui négocient le skew en sont les contreparties naturelles~; en
2012, les traders de la banque en vinrent à croire que ces contreparties connaissaient leurs positions.

\begin{omfigure}
\begin{tikzpicture}
\begin{axis}[width=11cm,height=4.8cm,
  xlabel={skew d'indice à la sortie (pb)},ylabel={P\&L (M USD)},
  tick label style={font=\small},label style={font=\small},
  xmin=-24,xmax=8,ymin=-8,ymax=6,grid=major,grid style={gray!25},
  legend columns=2,legend style={font=\footnotesize,at={(0.5,-0.3)},anchor=north,draw=none,fill=none,/tikz/every even column/.append style={column sep=8pt}},
  table/col sep=comma]
\addplot[omDef,very thick] table[x=skew,y=pnl]{figdata/markets-2/24-credit-indices-and-tranches/skewpnl.csv};
\addplot[omProp,thick,dashed] coordinates {(-24,1.53) (8,1.53)};
\addplot[only marks,mark=*,mark size=1.8pt,black,forget plot] coordinates {(-8,0) (0,3.55) (-20,-5.37)};
\legend{P\&L brut,coût aller-retour de négociation}
\end{axis}
\end{tikzpicture}

\omcaption{P\&L d'une transaction sur skew de 1~milliard de USD ouverte à un skew de $-8$ points de base
(achat de protection sur l'indice, vente de protection sur les constituants), en fonction du skew à la sortie,
les constituants restant inchangés. La transaction rapporte 3.55~millions de USD si le skew se referme et perd
5.37~millions de USD si le skew s'élargit d'abord à $-20$~; la ligne en pointillés donne le coût
d'entrée et de sortie des deux jambes. Illustratif~; données~: le problème du
week-end du chapitre.}\label{fig:m2:credit-indices-and-tranches:skewpnl}
\end{omfigure}

\section{Tranches}

\begin{definition}[Tranche, point d'attachement et point de détachement]\label{def:m2:credit-indices-and-tranches:tranche}
Une \emph{tranche}\index{tranche} d'un indice est un contrat qui ne couvre que les
pertes du portefeuille comprises entre deux niveaux, exprimées en fractions du notionnel de l'indice~:
les pertes commencent à la réduire au \emph{point d'attachement}\index{point d'attachement} $a$
et l'ont anéantie au \emph{point de détachement}\index{point de détachement} $d$. Pour une
perte $L$ du portefeuille, la tranche perd $\min\bigl(\max(L - a, 0),\, d - a\bigr)$, soit une fraction
$\min(\max(L-a,0), d-a)/(d-a)$ de son propre notionnel.
\end{definition}

Les \omterm{def:m2:credit-indices-and-tranches:tranche}{tranches} d'un indice se partagent ses pertes dans l'ordre
(\cref{fig:m2:credit-indices-and-tranches:stack})~: sur l'indice nord-américain
investment grade, la \omterm{def:m2:credit-indices-and-tranches:tranche}{tranche} equity, de 0 à 3~\%, absorbe les premières pertes et
rapporte la prime la plus élevée~; la \omterm{def:m2:credit-indices-and-tranches:tranche}{tranche} mezzanine, de 3 à 7~\%, n'est touchée qu'une fois la
\omterm{def:m2:credit-indices-and-tranches:tranche}{tranche} equity épuisée~; les \omterm{def:m2:credit-indices-and-tranches:tranche}{tranches} senior et super-senior situées au-dessus ne perdent qu'en cas de catastrophe (dans
l'exemple de ce chapitre, de 7 à 15~\% et de 15 à 100~\%). Avec un \omterm{def:m2:corporate-bonds:covenant}{taux de recouvrement} de 40~\%, chaque défaut d'un indice de 125 noms
coûte 0.48~\% du notionnel, si bien que la \omterm{def:m2:credit-indices-and-tranches:tranche}{tranche} equity est anéantie par sept défauts.
Les pertes espérées de l'ensemble des \omterm{def:m2:credit-indices-and-tranches:tranche}{tranches}, pondérées par leurs largeurs, s'additionnent pour donner la
perte espérée de l'ensemble du portefeuille, puisque ensemble elles le couvrent une seule fois. Ce qu'une
\omterm{def:m2:credit-indices-and-tranches:tranche}{tranche} ne partage pas avec l'indice, c'est sa dépendance à la manière dont les défauts se regroupent.

\begin{omfigure}
\centering
\begin{tikzpicture}[font=\small]
\fill[omDef!35] (0,0) rectangle (3.8,0.6);
\fill[omThm!30] (0,0.6) rectangle (3.8,1.4);
\fill[omProp!25] (0,1.4) rectangle (3.8,2.6);
\fill[gray!15] (0,2.6) rectangle (3.8,5.2);
\draw[thick] (0,0) rectangle (3.8,5.2);
\foreach \y in {0.6,1.4,2.6} \draw[thick] (0,\y) -- (3.8,\y);
\node at (1.9,0.3) {tranche equity 0--3~\%};
\node at (1.9,1.0) {tranche mezzanine 3--7~\%};
\node at (1.9,2.0) {tranche senior 7--15~\%};
\node at (1.9,3.9) {tranche super senior 15--100~\%};
\node[anchor=west,font=\footnotesize] at (4.0,0) {0~\%};
\node[anchor=west,font=\footnotesize] at (4.0,0.6) {3~\%};
\node[anchor=west,font=\footnotesize] at (4.0,1.4) {7~\%};
\node[anchor=west,font=\footnotesize] at (4.0,2.6) {15~\%};
\node[anchor=west,font=\footnotesize] at (4.0,5.2) {100~\%};
\draw[-{Stealth[length=2.5mm]},very thick,omDef] (-1.2,0) -- node[left,align=right,font=\footnotesize] {pertes du\\portefeuille\\remplissent la pile\\par le bas} (-1.2,2.4);
\node[align=left,font=\footnotesize,anchor=west] at (5.4,0.3) {première perte, prime la plus élevée~;\\anéantie par 7 défauts};
\node[align=left,font=\footnotesize,anchor=west] at (5.4,3.9) {ne perd que si plus de\\31 défauts surviennent};
\end{tikzpicture}

\omcaption{Les \omterm{def:m2:credit-indices-and-tranches:tranche}{tranches} de l'exemple de ce chapitre (hauteurs non à l'échelle). Les pertes du
portefeuille remplissent la pile par le bas~: la \omterm{def:m2:credit-indices-and-tranches:tranche}{tranche} equity absorbe les 3~\% premières pertes, la
mezzanine absorbe les 4~\% suivantes, et ainsi de suite~; avec un \omterm{def:m2:corporate-bonds:covenant}{taux de recouvrement} de 40~\%, chacun des 125 défauts coûte
0.48~\% du notionnel. Schéma de principe.}\label{fig:m2:credit-indices-and-tranches:stack}
\end{omfigure}

Le modèle standard de cette dépendance est le modèle gaussien à un facteur. Chaque nom
fait défaut avant l'horizon si $\sqrt\rho\,Z + \sqrt{1-\rho}\,\varepsilon_i <
N^{-1}(p)$, avec $Z$ un facteur commun, $\varepsilon_i$ des normales centrées réduites indépendantes, $p$
la probabilité de défaut à l'horizon et $\rho$ la corrélation entre les noms.
Sachant $Z$, les défauts sont indépendants et, dans un grand ensemble de noms, la fraction qui fait défaut est
leur probabilité conditionnelle, de sorte que le portefeuille perd
\[
L(Z) = (1 - R)\, N\!\left(\frac{N^{-1}(p) - \sqrt{\rho}\,Z}{\sqrt{1 - \rho}}\right),
\]
et la perte espérée d'une \omterm{def:m2:credit-indices-and-tranches:tranche}{tranche} est la moyenne de sa perte sur $Z$.

\begin{definition}[Corrélation de base]\label{def:m2:credit-indices-and-tranches:basecorr}
La \emph{corrélation de base}\index{corrélation de base} pour un \omterm{def:m2:credit-indices-and-tranches:tranche}{point de détachement} $d$ est la
corrélation $\rho$ pour laquelle le modèle à un facteur valorise la \omterm{def:m2:credit-indices-and-tranches:tranche}{tranche} allant de 0 à $d$ à
sa valeur de marché~; une \omterm{def:m2:credit-indices-and-tranches:tranche}{tranche} de $a$ à $d$ est alors valorisée comme la différence des
\omterm{def:m2:credit-indices-and-tranches:tranche}{tranches} de base $[0,d]$ et $[0,a]$, chacune à sa propre corrélation.
\end{definition}

La corrélation transfère le risque entre les \omterm{def:m2:credit-indices-and-tranches:tranche}{tranches}
(\cref{fig:m2:credit-indices-and-tranches:tranches}). Une corrélation faible rend les défauts
indépendants~: quelques défauts se produisent toujours, ce qui nuit à la \omterm{def:m2:credit-indices-and-tranches:tranche}{tranche} equity, et que de nombreux défauts ne surviennent presque jamais, ce qui
épargne la \omterm{def:m2:credit-indices-and-tranches:tranche}{tranche} senior. Une corrélation élevée fait que les défauts surviennent ensemble ou pas du tout, ce qui
épargne souvent la \omterm{def:m2:credit-indices-and-tranches:tranche}{tranche} equity et frappe parfois la \omterm{def:m2:credit-indices-and-tranches:tranche}{tranche} senior. La \omterm{def:m2:credit-indices-and-tranches:tranche}{tranche} equity est donc longue
en corrélation et la \omterm{def:m2:credit-indices-and-tranches:tranche}{tranche} senior courte. La mezzanine se situe entre les deux, et sa perte espérée
augmente d'abord puis diminue avec la corrélation, si bien qu'un seul prix de mezzanine peut correspondre à deux
corrélations~; la \omterm{def:m2:credit-indices-and-tranches:basecorr}{corrélation de base}, définie sur des \omterm{def:m2:credit-indices-and-tranches:tranche}{tranches} qui partent toutes de zéro, lève
cette ambiguïté, puisque la perte espérée d'une \omterm{def:m2:credit-indices-and-tranches:tranche}{tranche} de base diminue toujours lorsque la corrélation augmente.

\begin{example}[Pertes de tranches]\label{ex:m2:credit-indices-and-tranches:tranches}
Pour l'indice stylisé, la probabilité de défaut à cinq ans implicite dans le
\omterm{def:m2:credit-indices-and-tranches:intrinsic}{spread intrinsèque} est de 5.95~\%, et la perte espérée du portefeuille de 3.57~\%. Avec $\rho = 0.3$, les
pertes espérées représentent 59.6~\% de la \omterm{def:m2:credit-indices-and-tranches:tranche}{tranche} equity, 24.1~\% de la mezzanine, 7.83~\% de la
senior et 0.22~\% de la super senior~; pondérées par leurs largeurs, 3, 4, 8 et 85~\%, elles
s'ajoutent à 3.57~\%. La perte espérée de la \omterm{def:m2:credit-indices-and-tranches:tranche}{tranche} equity vaut 64.0~\% pour $\rho = 0.25$, et l'inversion
du modèle à cette valeur renvoie une \omterm{def:m2:credit-indices-and-tranches:basecorr}{corrélation de base} de 25~\%. Celle de la mezzanine est de 24.3~\%
à $\rho = 0.10$, 25.0~\% à 0.20 et 21.9~\% à 0.45~: elle vaut de nouveau 24.3~\% à 0.28.
\end{example}

\begin{omfigure}
\begin{tikzpicture}
\begin{axis}[width=11cm,height=5cm,
  xlabel={corrélation $\rho$},ylabel={perte espérée (\% de la tranche)},
  tick label style={font=\small},label style={font=\small},
  xmin=0.05,xmax=0.9,ymin=0.001,ymax=100,ymode=log,grid=major,grid style={gray!25},
  legend columns=4,legend style={font=\footnotesize,at={(0.5,-0.3)},anchor=north,draw=none,fill=none,/tikz/every even column/.append style={column sep=8pt}},
  table/col sep=comma]
\addplot[omDef,very thick] table[x=rho,y=equity]{figdata/markets-2/24-credit-indices-and-tranches/tranches.csv};
\addplot[omThm,very thick] table[x=rho,y=mezz]{figdata/markets-2/24-credit-indices-and-tranches/tranches.csv};
\addplot[omProp,very thick] table[x=rho,y=senior]{figdata/markets-2/24-credit-indices-and-tranches/tranches.csv};
\addplot[gray,very thick,dashed] table[x=rho,y=super]{figdata/markets-2/24-credit-indices-and-tranches/tranches.csv};
\legend{0--3~\%,3--7~\%,7--15~\%,15--100~\%}
\end{axis}
\end{tikzpicture}

\omcaption{Perte espérée à cinq ans de chaque \omterm{def:m2:credit-indices-and-tranches:tranche}{tranche}, en pourcentage de son notionnel et sur une
échelle logarithmique, en fonction de la corrélation du modèle gaussien à un facteur (probabilité de
défaut 5.95~\%, \omterm{def:m2:corporate-bonds:covenant}{taux de recouvrement} 40~\%). La perte espérée de la \omterm{def:m2:credit-indices-and-tranches:tranche}{tranche} equity décroît avec la corrélation
et celle de la super senior croît~; celle de la mezzanine croît puis décroît, et c'est aussi le cas, à
fortes corrélations, de celle de la \omterm{def:m2:credit-indices-and-tranches:tranche}{tranche} 7--15~\%. Illustratif~;
données~: le tutoriel du chapitre.}\label{fig:m2:credit-indices-and-tranches:tranches}
\end{omfigure}

\begin{omfigure}
\begin{tikzpicture}
\begin{axis}[width=11cm,height=5cm,
  xlabel={perte du portefeuille (\% du notionnel)},ylabel={probabilité},
  tick label style={font=\small},label style={font=\small},
  xmin=-0.3,xmax=19.5,ymin=0,ymax=0.28,grid=major,grid style={gray!25},
  yticklabel style={/pgf/number format/.cd,fixed,precision=2},
  legend columns=2,legend style={font=\footnotesize,at={(0.5,-0.3)},anchor=north,draw=none,fill=none,/tikz/every even column/.append style={column sep=8pt}},
  table/col sep=comma]
\addplot[omDef,very thick,mark=*,mark size=1.1pt] table[x=loss,y=r10]{figdata/markets-2/24-credit-indices-and-tranches/losses.csv};
\addplot[omThm,very thick,mark=*,mark size=1.1pt] table[x=loss,y=r40]{figdata/markets-2/24-credit-indices-and-tranches/losses.csv};
\addplot[gray,thick,dashed,forget plot] coordinates {(3,0) (3,0.28)};
\addplot[gray,thick,dashed,forget plot] coordinates {(7,0) (7,0.28)};
\addplot[gray,thick,dashed,forget plot] coordinates {(15,0) (15,0.28)};
\legend{$\rho = 0.1$,$\rho = 0.4$}
\end{axis}
\end{tikzpicture}

\omcaption{Distribution simulée de la perte à cinq ans du portefeuille de 125 noms, un
point par nombre de défauts (0.48~\% chacun), pour des corrélations de 0.1 et 0.4~; le dernier
point regroupe 40 défauts ou plus. Les traits pointillés marquent les
\omterm{def:m2:credit-indices-and-tranches:tranche}{points d'attachement} de 3, 7 et 15~\%. Une corrélation plus élevée accorde davantage de poids à l'absence de défaut et aux nombreux défauts. Illustratif~;
données~: le tutoriel du chapitre, 20\,000 tirages
à graine fixée.}\label{fig:m2:credit-indices-and-tranches:losses}
\end{omfigure}

Le modèle à un facteur est une convention de cotation, comme la formule de Black pour
les swaptions (chapitre 13), et non une description de la manière dont les défauts surviennent~; son échec
à ajuster toutes les \omterm{def:m2:credit-indices-and-tranches:tranche}{tranches} avec une même corrélation est le skew de corrélation~: un seul nombre ne peut décrire
comment les défauts se regroupent.

\section{Options sur indices}

\begin{definition}[Option sur indice de crédit]\label{def:m2:credit-indices-and-tranches:option}
Une \emph{option sur indice de crédit}\index{option sur indice de crédit} est une option permettant
d'entrer dans un contrat sur indice au spread d'exercice fixé à une date d'échéance~: une option
payeur donne le droit d'acheter la protection, une option receveur celui de la vendre.
\end{definition}

Les options sur indices se négocient sur les indices investment grade et high yield à cinq ans, de style
européen et réglées par livraison de la position sur l'indice, avec des échéances allant jusqu'à neuf
mois. Une option payeur vaut la peine d'être exercée si le spread de l'indice à l'échéance dépasse le strike~:
le détenteur achète la protection aux conditions du strike et reçoit la différence en upfront,
à peu près (spread moins strike) multiplié par l'annuité. Les payeurs sont le moyen par lequel les
investisseurs achètent de la protection contre un élargissement des \omterm{def:m2:corporate-bonds:spreads}{spreads de crédit} sans payer la prime courante,
et leur prix se calcule, en première approximation, avec la formule de Black appliquée au spread,
l'annuité jouant le rôle qu'elle joue pour les swaptions. Les détails, y compris la manière dont
les défauts survenant avant l'échéance sont traités et pourquoi le spread à terme doit les intégrer, relèvent
du livre~6.

\section{Tutoriel~: un indice face à ses constituants, et ses tranches}\label{tut:m2:credit-indices-and-tranches:tranche}

\begin{tutorial}
\textbf{Objectif.} Calculer le \omterm{def:m2:credit-indices-and-tranches:intrinsic}{spread intrinsèque} et le \omterm{def:m2:credit-indices-and-tranches:intrinsic}{skew d'indice} d'un indice stylisé à partir de ses 125
constituants, le P\&L d'une transaction de skew, et les pertes espérées des \omterm{def:m2:credit-indices-and-tranches:tranche}{tranches} dans le modèle
gaussien à un facteur, avec une \omterm{def:m2:credit-indices-and-tranches:basecorr}{corrélation de base}.
\textbf{État final~:} les \cref{fig:m2:credit-indices-and-tranches:constituents,fig:m2:credit-indices-and-tranches:skewpnl,fig:m2:credit-indices-and-tranches:tranches,fig:m2:credit-indices-and-tranches:losses},
et les \cref{ex:m2:credit-indices-and-tranches:intrinsic,ex:m2:credit-indices-and-tranches:tranches}
ainsi que les nombres du problème du week-end.
\begin{tutsteps}
\item \textbf{L'indice}~: upfront intrinsèque, spread pondéré par l'annuité, P\&L du skew.
\omcode{code/firm/tranche/firm_tranche.py}{17}{31}{Upfront intrinsèque, spread pondéré par l'annuité et P\&L d'une position de skew.}\label{lst:m2:credit-indices-and-tranches:index}
\item \textbf{Les \omterm{def:m2:credit-indices-and-tranches:tranche}{tranches}}~: perte de \omterm{def:m2:credit-indices-and-tranches:tranche}{tranche}, perte espérée en grand portefeuille, \omterm{def:m2:credit-indices-and-tranches:basecorr}{corrélation de base} et une
simulation sur un portefeuille fini.
\omcode{code/firm/tranche/firm_tranche.py}{36}{73}{Pertes espérées des tranches dans le modèle gaussien à un facteur.}\label{lst:m2:credit-indices-and-tranches:tranche}
\item \textbf{Exécuter} \texttt{index\_demo.tutorial()}, \texttt{index\_demo.skew\_trade()} et
\texttt{fig\_index.py}~; la conversion entre spreads et upfronts est celle du chapitre 23,
\texttt{firm.cds}.
\end{tutsteps}
\textbf{Pour aller plus loin.} Faire bondir un constituant de 100 à 1\,000 points de base et
observer comment le \omterm{def:m2:credit-indices-and-tranches:intrinsic}{spread intrinsèque} et les \omterm{def:m2:credit-indices-and-tranches:tranche}{tranches} évoluent~; puis remplacer le grand portefeuille par
la simulation finie à 125 noms et mesurer l'erreur sur la \omterm{def:m2:credit-indices-and-tranches:tranche}{tranche} equity.
\end{tutorial}

\section{Construction~: analytique de l'indice et des tranches}\label{bld:m2:credit-indices-and-tranches:tranche}

\begin{build}
\bfield{Rôle} La firme miniature valorise ses positions sur indices par rapport à leurs
constituants, surveille le skew et évalue le risque des \omterm{def:m2:credit-indices-and-tranches:tranche}{tranches} pour ses modèles de crédit.
\bfield{Interface} \texttt{intrinsic\_upfront(upfronts, weights)}~;
\texttt{annuity\_weighted\_spread(spreads, annuities)}~; \texttt{skew\_pnl(notional,
upfront\_entry, upfront\_exit)}~; \texttt{tranche\_loss(pool\_loss, attach, detach)}~;
\texttt{expected\_tranche\_loss(attach, detach, p, rho, recovery)}~;
\texttt{base\_correlation(detach, target, p, recovery)}~; \texttt{simulate\_pool(names, p,
rho, recovery, trials, seed)}. \bfield{Règles} Pondérations égales sauf indication contraire~; upfronts par unité de notionnel issus de
\texttt{firm.cds}~; grand portefeuille homogène avec une seule probabilité de défaut et un seul
\omterm{def:m2:corporate-bonds:covenant}{taux de recouvrement}~; pertes espérées par la règle des trapèzes sur le facteur dans $[-8, 8]$.
\bfield{Tests d'acceptation} \texttt{code/firm/tranche/tests/}~: pertes de \omterm{def:m2:credit-indices-and-tranches:tranche}{tranche} bornées à
$[0, 1]$~; les pertes espérées des \omterm{def:m2:credit-indices-and-tranches:tranche}{tranches} se somment à celles du portefeuille~; la \omterm{def:m2:credit-indices-and-tranches:tranche}{tranche} equity baisse et la \omterm{def:m2:credit-indices-and-tranches:tranche}{tranche} senior monte
avec la corrélation~; la \omterm{def:m2:credit-indices-and-tranches:basecorr}{corrélation de base} fait l'aller-retour~; un portefeuille fini de 125 noms concorde avec le
grand portefeuille.
\bfield{Pour aller plus loin} Probabilités de défaut hétérogènes avec la récurrence pour la
distribution des pertes du portefeuille fini~; \omterm{def:m2:corporate-bonds:covenant}{taux de recouvrement} stochastique~; bootstrap des \omterm{def:m2:credit-indices-and-tranches:basecorr}{corrélations de base}
à partir des cotations de \omterm{def:m2:credit-indices-and-tranches:tranche}{tranches}~; pricing d'options sur indices (livre~6).
\end{build}

\begin{omsources}
\item SEC, Order Instituting Cease-and-Desist Proceedings, In the Matter of JPMorgan Chase \& Co., Release No.~34-70458, 19 septembre 2013. \item JPMorgan Chase \& Co., Report of the Management Task Force Regarding 2012 CIO
Losses, 16 janvier 2013.
\item US Senate, Permanent Subcommittee on Investigations, audition «~JPMorgan Chase
whale trades: a case history of derivatives risks and abuses~», 15 mars 2013.
\item ICE, Markit CDX.NA.IG contract specifications~; CDS index options clearing.
\item Bloomberg SEF, CDX.NA.IG five-year contract submission to the CFTC, novembre 2013.
\end{omsources}

\section{Exercices}

\begin{exercise}[$\star$]\label{exo:m2:credit-indices-and-tranches:1}
Un investisseur a vendu 1~milliard de dollars US de protection sur un indice à 125 noms. Un des noms fait défaut et son
\omterm{def:m2:credit-default-swaps:final}{prix final d'enchère} est de 40. Quel montant faut-il payer, et quel est le notionnel de l'indice ensuite~?
\end{exercise}

\begin{exercise}[$\star$]\label{exo:m2:credit-indices-and-tranches:2}
Quelle fraction de la \omterm{def:m2:credit-indices-and-tranches:tranche}{tranche} 3--7~\% est perdue si les pertes du portefeuille sont de 2~\%, 5~\% et 9~\%~?
\end{exercise}

\begin{exercise}[$\star$]\label{exo:m2:credit-indices-and-tranches:3}
Pourquoi une nouvelle \omterm{def:m2:credit-indices-and-tranches:index}{série d'indice} concentre-t-elle l'essentiel des échanges, et que devient l'ancienne~?
\end{exercise}

\begin{exercise}[$\star\star$]\label{exo:m2:credit-indices-and-tranches:4}
Un indice compte deux noms, à 50 et 250 points de base. Son \omterm{def:m2:credit-indices-and-tranches:intrinsic}{spread intrinsèque} vaut-il 150~? Utiliser
les conventions du chapitre.
\end{exercise}

\begin{exercise}[$\star\star$]\label{exo:m2:credit-indices-and-tranches:5}
Pourquoi un \omterm{def:m2:credit-default-swaps:cds}{vendeur de protection} sur la \omterm{def:m2:credit-indices-and-tranches:tranche}{tranche} equity est-il long de corrélation~?
\end{exercise}

\begin{exercise}[$\star\star$]\label{exo:m2:credit-indices-and-tranches:6}
Un investisseur achète une option payeuse à trois mois sur l'indice, dont le prix d'exercice est de 80 points de base.
Quand l'option est-elle exercée, et que reçoit l'investisseur~?
\end{exercise}

\begin{exercise}[$\star\star\star$]\label{exo:m2:credit-indices-and-tranches:7}
\emph{Programmation.} Avec \texttt{expected\_tranche\_loss}, vérifier que les pertes espérées des quatre \omterm{def:m2:credit-indices-and-tranches:tranche}{tranches}
s'additionnent à celles du portefeuille pour $\rho = 0.3$, et trouver la seconde
valeur de corrélation pour laquelle la perte espérée de la \omterm{def:m2:credit-indices-and-tranches:tranche}{tranche} mezzanine égale sa valeur en $\rho = 0.1$.
\end{exercise}

\begin{exercise}[$\star\star\star$]\label{exo:m2:credit-indices-and-tranches:8}
\emph{Repérer l'erreur.} «~L'indice est la moyenne de ses 125 spreads. Il se négocie 8 points de base
sous cette moyenne, donc acheter la protection sur l'indice et vendre les noms~: un profit sans
risque.» Corriger cette affirmation.
\end{exercise}

\section{Problème~: le skew trade}

\begin{problem}[{Problème du week-end --- trading d'un indice contre ses constituants}]\label{pb:m2:credit-indices-and-tranches:1}
Un fonds de relative value suit l'indice stylisé de
l'\cref{ex:m2:credit-indices-and-tranches:intrinsic}~: 125 noms, \omterm{def:m2:credit-indices-and-tranches:intrinsic}{spread intrinsèque} de 73.6
points de base, tandis qu'un gros \omterm{def:m2:credit-default-swaps:cds}{vendeur de protection} sur l'indice l'a poussé à 65.6.
Le fonds achète 1~milliard de dollars US de protection sur l'indice et vend 8~millions de dollars US de protection
sur chaque constituant. Négocier un seul nom coûte une demi-fourchette de 1.5 points de base
de spread courant, l'indice 0.25~; taux sans risque constant de 4~\%, \omterm{def:m2:corporate-bonds:covenant}{taux de recouvrement} de 40~\%.

\medskip\noindent\textbf{Partie I --- L'indice et ses constituants.}
\begin{enumerate}
\item Quelles sont la moyenne simple et la médiane des spreads des constituants~?
\item Quel est l'upfront intrinsèque, et quel est le \omterm{def:m2:credit-indices-and-tranches:intrinsic}{spread intrinsèque}~?
\item Pourquoi le \omterm{def:m2:credit-indices-and-tranches:intrinsic}{spread intrinsèque} est-il inférieur à la moyenne simple~?
\item Comparer le \omterm{def:m2:credit-indices-and-tranches:intrinsic}{spread intrinsèque} à la moyenne des spreads pondérée par les annuités.
\item Quel est le skew~?
\end{enumerate}

\medskip\noindent\textbf{Partie II --- Le trade.}
\begin{enumerate}[resume]
\item Quel upfront le fonds reçoit-il sur l'indice, et quel upfront verse-t-il sur les constituants~?
\item Quel est le flux de trésorerie net à l'entrée, et que représente-t-il~?
\item Un constituant fait défaut, avec un \omterm{def:m2:credit-default-swaps:final}{prix final d'enchère} de 40. Que paie le fonds
et que reçoit-il~?
\item Quel est le P\&L si le skew se referme, les constituants restant inchangés~?
\item Quel est le P\&L si le skew s'élargit d'abord à $-20$ points de base~?
\end{enumerate}

\medskip\noindent\textbf{Partie III --- Coûts et risques.}
\begin{enumerate}[resume]
\item Quel est le coût d'entrée dans les deux jambes, et celui d'entrée puis de sortie~?
\item Quel est le P\&L net de coûts si le skew se referme, et de combien le skew doit-il
évoluer pour que le trade couvre ses coûts~?
\item Pourquoi le skew peut-il persister ou s'élargir~?
\item Que coûte un élargissement au fonds avant son terme, au-delà de la perte latente~?
\item Pourquoi un très gros \omterm{def:m2:credit-default-swaps:cds}{vendeur de protection} sur indices est-il exposé aux traders de skew~?
\end{enumerate}

\medskip\noindent\textbf{Partie IV --- Jugement.}
\begin{enumerate}[resume]
\item Pourquoi le vendeur ne pourrait-il pas simplement dénouer une position de 10 à 15 jours de
volume du marché~?
\item Comment dimensionner une transaction sur skew~?
\item Que devient la transaction au prochain roll~?
\item Énoncer le \emph{résultat nommé}~: le P\&L de la transaction sur skew lorsque le skew se referme,
brut et net de coûts, et sa perte si le skew s'élargit d'abord.
\item En une phrase~: pourquoi le skew n'est-il pas un gain sans risque~?
\end{enumerate}
\end{problem}

\section{Questions d'entretien}

\begin{interviewq}[$\star$ \iqroles{trader, chercheur}]\label{iq:m2:credit-indices-and-tranches:1}
Qu'est-ce que le CDX investment grade, et comment se produit son roll~?
\end{interviewq}

\begin{interviewq}[$\star$ \iqroles{trader}]\label{iq:m2:credit-indices-and-tranches:2}
Qu'est-ce que le \omterm{def:m2:credit-indices-and-tranches:intrinsic}{spread intrinsèque} d'un indice, et pourquoi l'indice peut-il s'écarter de celui-ci~?
\end{interviewq}

\begin{interviewq}[$\star\star$ \iqroles{chercheur}]\label{iq:m2:credit-indices-and-tranches:3}
Dérivez la perte d'un grand panier dans le modèle gaussien à un facteur.
\end{interviewq}

\begin{interviewq}[$\star\star$ \iqroles{chercheur, trader}]\label{iq:m2:credit-indices-and-tranches:4}
Quelles \omterm{def:m2:credit-indices-and-tranches:tranche}{tranches} sont longues en corrélation, et pourquoi~? Pourquoi utiliser la \omterm{def:m2:credit-indices-and-tranches:basecorr}{corrélation de base}
plutôt qu'une corrélation par \omterm{def:m2:credit-indices-and-tranches:tranche}{tranche}~?
\end{interviewq}

\begin{interviewq}[$\star\star$ \iqroles{risque, trader}]\label{iq:m2:credit-indices-and-tranches:5}
Qu'est-ce qui a mal tourné dans les pertes de 2012 du Chief Investment Office de JPMorgan~?
\end{interviewq}

\begin{interviewq}[$\star\star\star$ \iqroles{développeur}]\label{iq:m2:credit-indices-and-tranches:6}
Concevez un service qui calcule, en temps réel, le \omterm{def:m2:credit-indices-and-tranches:intrinsic}{spread intrinsèque} et le \omterm{def:m2:credit-indices-and-tranches:intrinsic}{skew d'indice} de chaque
\omterm{def:m2:credit-indices-and-tranches:index}{série d'indice} à partir des cotations sur noms individuels.
\end{interviewq}
